\section{Open Weights Go Architectural}
\sectionkicker{OPEN WEIGHTS}
\begin{wideflow}
{\Large\bfseries 公開の重みは構造で競う}\par
\smallskip
{\color{SurveyMuted}大きなMoEの公開と予告が並んだ週だが、動作規模の小ささも利用条件もモデルごとに確かめるべきで、ひと括りの効率物語にはしない。}\par
\end{wideflow}

8月26日に報じられたZ.aiのGLM-5.3-Flashは、8月20日からOpenRouter上で無償のステルス preview として現れていたOx Alphaがその正体だったという時系列で語られている\cite{glm53flash}。320BのMoEに対して動作時約18Bという規模感はこのモデル固有の数値として伝えられ、比率にして約5.6\%に相当するが、他のモデルに広げられる性質ではない。重みのMIT表示は重み置き場での確認を待つ段階であり、1Mトークンの入力や131Kの出力、疎注意と線形注意の併用、30Tトークンの学習といった仕様も報道経由の数値として扱う。GDPval-AAでの上位やAutomationBenchでの2位、10倍の費用主張はいずれも提供側の発表として受け止め、独立の再現を待つ。

\begin{themeoverview}[三つの公開と一つの読み方]
\begin{tabularx}{\linewidth}{@{}>{\bfseries}p{0.25\linewidth}X@{}}
GLM-5.3-Flash & 320B中約18Bの動作という固有の規模とOx Alphaからの時系列を報道経由で示す\cite{glm53flash}。\\
Qwen3.8-Flash-Next & 125B中約6Bの動作という固有の規模と次世代予告の位置づけを報道経由で示す\cite{qwen38flash}。\\
Hy4 preview & 770B中約49Bの動作という固有の規模と preview の注意書きを報道経由で示す\cite{hy4preview}。\\
\end{tabularx}
\end{themeoverview}

8月26日に重みが公開されたQwen3.8-Flash-Nextは、125BのMoEに対して動作時約6Bに51Bのn-gram表を伴う構成として報じられ、比率にして約4.8\%という低さはあくまでこのモデル固有の話である\cite{qwen38flash}。カード上の実験名が次世代の予告を示すことから旗艦の前触れとして読めるが、Qwen 4そのものが出たわけではなく、期日や重み、カードは確認されていない。利用条件の条文も一次資料での確認を待つ。vLLMやSGLangへの対応 pull request は未結合の方向付けに留まり、符号や符号審査の数値や秒間540トークンという提供側数値は再現前の主張として区別する。

8月28日に報じられたTencentのHy4は、符号処理や研究、金融分析向けの preview として重み置き場経由で出たという位置づけである\cite{hy4preview}。770BのMoEに対して要求あたり約49Bが動くという規模感はこの preview 固有の数値として報じられ、比率にして約6.4\%に相当する。CodeBuddyやWorkBuddyとの連携は予定として語られる。提供側自身が初期版は考え過ぎや検証過多になり得ると述べている点は重要で、読者はこの注意書きごと受け止める必要がある。順位や測定値の裏づけは今回の資料にはなく、窓末の公表ゆえ独立の検証の足跡もまだ薄い。

\begin{claimboundary}[未解決の境界]
三つの動作比率を一つの共通効率として束ねない。利用条件もモデルごとに一次確認を待ち、共通の許諾や共通の待避・記憶機構として語らない。提供側の測定値は再現まで提供側の主張であり、Qwen 4の旗艦を発売済みとして扱わない。報道経由の仕様や価格は重み置き場や一次発表での確認を欠く限り事実に格上げしない。
\end{claimboundary}

この三件の周りでは読者の関心も集まり、公式の発信に複数の独立した言及が重なる様子が観察記録に残る\cite{w35community}。ただし観察記録は関心の方向を示すだけで、仕様や測定値、日付、利用条件、公開事実の根拠にはならない。一方、IBMが8月25日に示したGranite 4.2は3B、8B、30Bの高密度モデル群で、8Bと30Bに複数段の強化学習による道具型訓練を施した点が特徴として報じられる\cite{granite42}。Apache 2.0の表示は重み置き場での確認を待ち、学習資料の7割近くが符号寄りという偏りは読みの前提として残す。測定値は提供側の数値であり、符号外への転移は確かめられていない。

\begin{technicalnote}[公開情報の境界]
X上の反応は関心の方向を示す背景にとどめる。技術仕様、性能、互換性は公式資料や重み置き場の記録から読み取り、反応だけから新しい公開や性能を推論しない\cite{w35community}。
\end{technicalnote}
